\section{Proof of the functional and joint limits}\label{sec:process}
\begin{proof}[Proof of Theorem~\ref{thm:main}]
Fix $\eta\in C_c^\infty$ equal to one near $\operatorname{supp}f_0$.
Initially every particle lies in that support, so for each $q\in\cQ$,
$\pi_0^\eps(q)=\pi_0^\eps(\eta q)$ almost surely, including equality of
expectations. Input F therefore gives the initial Gaussian vector $Z_0$
with covariance \eqref{eq:initial-cov}. This is an exact compact-test
substitution at time zero, not an unbounded positive-time observation.
Lemma~\ref{lem:vacancy} adds its deterministic vector $a$ for Boltzmann
centering. Since $r_q(t)$ is a fixed finite linear combination of the basis
\eqref{eq:basis}, all of its initial fluctuations follow jointly from
this one vector.

Subtract \eqref{eq:Bol-Q} from \eqref{eq:Q-identity} and divide by $\eps$:
\begin{equation}\label{eq:exact-path}
 Y_t^\eps(q)=Y_0^\eps(r_q(t))
       +d(q)\cH_\eps(t)+2e(q)\int_0^t\cH_\eps(s)\dd s.
\end{equation}
For the pasted law use \eqref{eq:L2-LLN}; for the true law use
\eqref{eq:true-LLN}. The event term converges uniformly in probability
to its continuous deterministic limit. The initial-vector map is continuous
into the space of continuous finite-coefficient polynomial paths. Ordinary
finite-dimensional Slutsky gives the joint $D([0,T],\R^{13})$, $J_1$
limit \eqref{eq:main-limit}. Indeed \eqref{eq:exact-path} provides a
uniform-in-time residual tending to zero. No independence between the
initial vector and the speed sum is required, since the latter limit
is deterministic.

Taking expectations in the same identity, using the vacancy bias and the
separate mean convergence of both speed sums, proves
\eqref{eq:mean-limit}. Uniformity follows from the fixed polynomial
coefficients and the uniform expectation limit. Subtracting the convergent
means proves the exactly centered limit. Formula \eqref{eq:limit-cov}
is then the covariance of the Gaussian linear initial field.

For the pasted prelimit covariance, let
$Z_{0,\eps}(q)=\zeta_0^\eps(q)$ and
$R_{q,\eps}(t)=d(q)[\cH_\eps^p(t)-m_\eps(t)]
+2e(q)\int_0^t[\cH_\eps^p(s)-m_\eps(s)]\dd s$.
The centered path identity reads
\[
 \zeta_t^{\eps,p}(q)=Z_{0,\eps}(r_q(t))+R_{q,\eps}(t).
\]
Proposition~\ref{prop:LLN} implies
$\E\sup_t|R_{q,\eps}(t)|^2\to0$.
The initial covariance matrices converge to \eqref{eq:initial-cov} by
the initial static factorial bounds, and are bounded. Finite coefficients
in $r_q(t)$ and Cauchy--Schwarz make every covariance involving a residual
uniformly small in $t,s$. The remaining initial covariance converges
uniformly to \eqref{eq:limit-cov}. This argument transfers no squared
bad-event error to the true law.

It remains to justify the product with the strong-dual field. Let
$\mathfrak E=D([0,T],\cS'_\beta)$ have input F's generalized $J_1$
topology, and let $\mathcal Z_\eps$ be its exactly centered field.
Evaluation at zero against a fixed finite Schwartz list is continuous:
every allowed time change fixes zero. With the basis $q_1,\ldots,q_{13}$,
\[
 Y_0^\eps(q_i)=\mathcal Z_\eps(0)(\eta q_i)+a_\eps(q_i).
\]
Define the continuous map from $\mathfrak E$ to
$\mathfrak E\times D([0,T],\R^{13})$ which retains its first coordinate
and constructs its second by \eqref{eq:main-limit}, replacing
$Z_0(q_i)$ by $z(0)(\eta q_i)$. Input F and continuous mapping give its
weak limit. The actual second coordinate differs uniformly in probability
from this construction by \eqref{eq:exact-path} and
$a_\eps\to a$.

For completeness, this perturbation step can be proved without assuming a
general nonmetrizable Slutsky theorem. Given a bounded continuous function
$g$ on the product and a compact $K\subset\mathfrak E$ carrying all but
arbitrarily small probability of the tight field laws, the constructed
graph over $K$ is compact. At each graph point continuity gives a field
neighborhood and a radius in the metric $J_1$ second factor where $g$
changes by less than a chosen $\delta$. Shrink the field neighborhood
using continuity of the constructed second coordinate. A finite subcover
of $K$ gives a common positive radius. If the field lies in $K$ and the
second-coordinate perturbation is smaller than this radius, the two
$g$ values differ by at most $2\delta$. Its exceptional probability
tends to zero. Outside $K$, use $2\|g\|_\infty$; then let the compact
exception and $\delta$ tend to zero. This proves convergence of all
bounded continuous product tests, as required. The joint conclusion
follows, with $Z_0(q)=G_0(\eta q)$.
\end{proof}

\begin{proof}[Proof of Corollary~\ref{cor:three}]
The transported tests are
$r_E=|v|^2$, $r_V=x\cdot v+t|v|^2$, and $r_X=|x+tv|^2$.
Their initial-coefficient matrix is $A(t)$ in \eqref{eq:three-limit}.
For virial $d=1,e=0$, and for spread $d=0,e=1$; energy has both zero.
Theorem~\ref{thm:main} gives all assertions. Fubini gives
$2\int_0^tH(s)\dd s=\int_0^t(t-s)J_f(s)\dd s$.
\end{proof}

The joint Gaussian cross covariance, for $q\in\cQ$ and
$\phi\in\cS$, is
\begin{equation}\label{eq:cross-cov}
 \operatorname{Cov}(G_s(\phi),Y_q(t))
           =\int f_0(z)[U(0,s)\phi](z)r_q(t,z)\dd z.
\end{equation}
Replace $r_q$ by $\eta r_q$ on the compact initial support and use
independence of the initial Gaussian field and its noise in input F.
This proves \eqref{eq:cross-cov} using only the existing
compact/Schwartz interpretation of $U$. The phase-space field in the
joint theorem retains exact microscopic centering; no generic
Boltzmann recentering of that field is implicit.

The collisional drift is nonzero for the stated smooth nonzero datum.
At some open set of positions the initial velocity density is positive
on an open velocity set, so
$\int |v-w|^2f_0(x,v)f_0(x,w)\dd x\dd v\dd w>0$.
Continuity on a compact subregion preserves a positive lower bound for
small positive times. Thus $H(t)>0$ there. The Gaussian part, by contrast,
comes entirely from the initial field because all transported tests in
$\cQ$ are co-located collision invariants.
